% !TEX program = platex
% 数学ノート集・共通テンプレート v2（積分論2026年度第1部の版面を参考）
% A5判。章内で定義・定理・例・注意等の番号を共通化。
% コンパイル： platex -interaction=nonstopmode math_note_template_v2.tex (2回)
%              dvipdfmx math_note_template_v2.dvi
\documentclass[a5paper,10pt,onecolumn]{jsarticle}
\usepackage[dvipdfmx,top=22mm,bottom=20mm,inner=19mm,outer=18mm,headheight=13pt,headsep=6mm]{geometry}
\usepackage{amsmath,amssymb,amsthm,mathtools}
\usepackage{enumitem,fancyhdr,titlesec}
\usepackage[dvipdfmx,bookmarks=true,colorlinks=false]{hyperref}
\usepackage{pxjahyper}

% ----- 差し替える文書情報 -----
\newcommand{\NoteTitle}{計数法と組合せ論的反転}
\newcommand{\NoteSeries}{数学ノート集}

% ----- 文字・組版 -----
% 参考PDFの質素な明朝＋Computer Modernの印象を引き継ぎながら，
% A5判で窮屈にならないよう本文の行間を少し確保する．
\linespread{1.13}
\setlength{\parindent}{1zw}
\setlength{\parskip}{0pt}
\setlength{\abovedisplayskip}{6pt plus 2pt minus 2pt}
\setlength{\belowdisplayskip}{6pt plus 2pt minus 2pt}
\setlength{\abovedisplayshortskip}{3pt plus 2pt minus 1pt}
\setlength{\belowdisplayshortskip}{4pt plus 2pt minus 1pt}
\setlist[enumerate]{topsep=3pt,itemsep=1pt,parsep=0pt,leftmargin=2.4zw}
\titleformat{\section}{\normalfont\normalsize\bfseries}{\thesection}{0.8zw}{}
\titlespacing*{\section}{0pt}{2.0ex plus .4ex}{.9ex}
\titleformat{\subsection}{\normalfont\normalsize\bfseries}{\thesubsection}{0.7zw}{}
\titlespacing*{\subsection}{0pt}{1.5ex plus .3ex}{.6ex}
\raggedbottom

% ----- 柱とページ番号：ページ上部、罫線なし -----
\fancypagestyle{mathnote}{%
  \fancyhf{}%
  \fancyhead[L]{\footnotesize \NoteSeries\ （\NoteTitle）}%
  \fancyhead[R]{\footnotesize \thepage}%
  \renewcommand{\headrulewidth}{0pt}%
  \renewcommand{\footrulewidth}{0pt}%
}
\pagestyle{mathnote}

% ----- 節内通し番号の定義・定理群 -----
\newtheoremstyle{mathnotetheorem}
 {4.5pt}{4.5pt}{\normalfont}{}%
 {\bfseries}{}{0.65zw}%
 {\thmname{#1}\thmnumber{ #2}\thmnote{（#3）}}
\theoremstyle{mathnotetheorem}
\newtheorem{theorem}{定理}[section]
\newtheorem{definition}[theorem]{定義}
\newtheorem{lemma}[theorem]{補題}
\newtheorem{proposition}[theorem]{命題}
\newtheorem{corollary}[theorem]{系}
\newtheorem{example}[theorem]{例}
\newtheorem{remark}[theorem]{注意}
\renewcommand{\proofname}{証明}
\renewcommand{\qedsymbol}{$\square$}
\makeatletter
\renewenvironment{proof}[1][\proofname]{%
  \par\pushQED{\qed}\normalfont\topsep4\p@\@plus2\p@\relax
  \trivlist\item[\hskip\labelsep\bfseries #1]\ignorespaces
}{\popQED\endtrivlist\@endpefalse}
\makeatother

% ----- 使うノートで適宜拡張する記号 -----
\newcommand{\R}{\mathbb{R}}
\newcommand{\norm}[1]{\lVert #1\rVert}
\newcommand{\abs}[1]{\lvert #1\rvert}

\begin{document}
\section*{読了後にできること}
本稿を読むための前提は，有限集合の写像，包除原理，Möbius反転．
本稿の具体的な到達課題は，包除によって全射の個数を求め，逆像の分類を具体的に実行する．

\section{計数の基本原理}
\begin{proposition}[加法原理と乗法原理]
互いに素な有限集合 $A,B$ に対して
$|A\cup B|=|A|+|B|$ が成立する．
また任意の有限集合 $A,B$ に対し
$|A\times B|=|A||B|$ である．
\end{proposition}
\begin{proof}
前者は互いに素な集合の元を別々に数えればよい．
後者は $A\times B$ が相互に素な集合
$\{a\}\times B$（$a\in A$）の和であり，
各集合が $B$ と全単射であることから従う．
\end{proof}
\begin{proposition}[写像の数]
$|A|=n$，$|B|=m$ なら $A\to B$ の写像は $m^n$ 個ある．
$n\le m$ の場合，そのうち単射は
$m(m-1)\cdots(m-n+1)$ 個である．
\end{proposition}
\begin{proof}
$A=\{a_1,\ldots,a_n\}$ と書くと，
写像は順序付き組 $(f(a_1),\ldots,f(a_n))$ と
一対一に対応するので $m^n$ 個である．
単射なら第 $k$ 番目の値の選択肢は
既に使った $k-1$ 個を除く $m-k+1$ 個だから，
乗法原理より表示式を得る．
\end{proof}

\begin{definition}
整数 $n\ge0$，$0\le k\le n$ に対し
$ {n\choose k}=n!/(k!(n-k)!)$ と定める．
範囲外の $k$ では $ {n\choose k}=0$ と約束する．
\end{definition}
\begin{lemma}\label{lem:choose}
$n$ 元集合の $k$ 元部分集合はちょうど
$ {n\choose k}$ 個ある．
\end{lemma}
\begin{proof}
相異なる $k$ 個を順序付きで選ぶ方法は
$n!/(n-k)!$ 通りである．
各 $k$ 元部分集合は $k!$ 通りの順序を持つので割り算する．
$k=0$ には空集合がただ一つ対応する．
\end{proof}
\begin{theorem}[PascalとVandermondeの恒等式]
$n\ge1$，$1\le k\le n$ なら
\[
 {n\choose k}={n-1\choose k-1}+{n-1\choose k}.
\]
また $r,s,k\ge0$ なら
\[
 {r+s\choose k}=\sum_{j=0}^k{r\choose j}{s\choose k-j}.
\]
\end{theorem}
\begin{proof}
第一式では一つの元を固定し，
それを選ぶ場合と選ばない場合を数える．
第二式では互いに素な $r$ 元集合と $s$ 元集合の和から
$k$ 元を選ぶ際，前者から選ぶ元の数を $j$ として場合分けする．
各場合は互いに素なので数を加えてよい．
\end{proof}
\begin{theorem}[多項定理]
$n_i\ge0$，$\sum_{i=1}^r n_i=n$ とする．
$n$ 元集合を番号付きの $r$ 個の箱へ
それぞれ $n_i$ 個ずつ割り当てる方法は
$n!/(n_1!\cdots n_r!)$ 通りである．従って
\[
 (x_1+\cdots+x_r)^n
 =\sum_{n_1+\cdots+n_r=n}
 {n!\over n_1!\cdots n_r!}x_1^{n_1}\cdots x_r^{n_r}.
\]
\end{theorem}
\begin{proof}
最初の箱へ $ {n\choose n_1}$ 通り，
二番目へ $ {n-n_1\choose n_2}$ 通り，順に選ぶと
その積は表示した階乗の比に約分される．
展開された項は $n$ 個の因子から選択した項の列に対応し，
各指数を固定した列の個数が前半の値に等しい．
\end{proof}
\section{包除原理と全射}
\begin{theorem}[包除原理]\label{thm:ie}
有限集合 $X$ とその部分集合 $A_1,\ldots,A_r$ について
\[
 \left|X\setminus\bigcup_{i=1}^rA_i\right|
 =\sum_{J\subset\{1,\ldots,r\}}(-1)^{|J|}
 \left|\bigcap_{j\in J}A_j\right|.
\]
空な添字集合についての共通部分は $X$ とする．
\end{theorem}
\begin{proof}
$x\in X$ に対し $I_x=\{i:x\in A_i\}$ と置く．
右辺で $x$ が数えられる重みの合計は
$\sum_{J\subset I_x}(-1)^{|J|}=(1-1)^{|I_x|}$ である．
これは $I_x=\varnothing$ なら1，
それ以外なら0となり，左辺と一致する．
全ての $x\in X$ について加えればよい．
\end{proof}
\begin{corollary}[全射の個数]\label{cor:surj}
$n$ 元集合から $r$ 元集合への全射は
\[
 \sum_{j=0}^r(-1)^j{r\choose j}(r-j)^n
\]
個である．ただし $0^0=1$ とする．
\end{corollary}
\begin{proof}
$A_i$ を値域の第 $i$ 元を取らない写像の集合とする．
指定した $j$ 個の値を避ける写像は $(r-j)^n$ 個であり，
全射は全ての $A_i$ を避ける写像に等しい．
定理\ref{thm:ie}を適用すればよい．
\end{proof}

\begin{definition}
$n$ 元集合を空でない $k$ 個の番号なしの部分集合へ
分割する方法の数を第二種Stirling数 $S(n,k)$ とする．
$S(0,0)=1$，$S(n,0)=0$（$n>0$）とする．
\end{definition}
\begin{proposition}[Stirling数の漸化式]
$n,k\ge1$ に対し
\[
 S(n,k)=S(n-1,k-1)+kS(n-1,k)
\]
が成り立つ．さらに
\[
 k!S(n,k)=\sum_{j=0}^{k}(-1)^j{k\choose j}(k-j)^n.
\]
\end{proposition}
\begin{proof}
最後の元が一元のみのブロックを作る場合には，
残る $n-1$ 元を $k-1$ ブロックに分割する．
そうでない場合には $n-1$ 元を $k$ ブロックに分割し，
最後の元をいずれかのブロックへ加える．
この二つの場合は互いに素なので漸化式を得る．
各分割の $k$ ブロックに番号を付ける方法は $k!$ 通りである．
番号付き分割は全射と一対一に対応するので，
第二式は系\ref{cor:surj}から従う．
\end{proof}
\begin{example}[固定点を持たない置換]
$n$ 元集合上の固定点を持たない置換の個数を $D_n$ とすると
\[
 D_n=n!\sum_{j=0}^n{(-1)^j\over j!}.
\]
\end{example}
\begin{proof}
全置換 $n!$ 個を考え，
$A_i$ を第 $i$ の元を固定する置換の集合とする．
$j$ 個の指定された点を固定する置換は $(n-j)!$ 個ある．
包除原理によって
\[
 D_n=\sum_{j=0}^n(-1)^j{n\choose j}(n-j)!
\]
を得る．二項係数を階乗で表せば結論となる．
\end{proof}
\section{二項反転}
\begin{theorem}[二項反転公式]\label{thm:binomialinv}
可換環 $R$ の元の二数列 $(f_n),(g_n)$ について
\[
 f_n=\sum_{k=0}^n{n\choose k}g_k
\quad(n\ge0)
\]
であることと
\[
 g_n=\sum_{k=0}^n(-1)^{n-k}{n\choose k}f_k
\quad(n\ge0)
\]
であることは同値である．
\end{theorem}
\begin{proof}
第一式を第二式の右辺に代入すると，
$g_j$ の係数は
\[
 \begin{aligned}
 &\sum_{k=j}^n(-1)^{n-k}{n\choose k}{k\choose j}\\
 &\quad={n\choose j}
  \sum_{k=j}^n(-1)^{n-k}{n-j\choose k-j}.
 \end{aligned}
\]
ここで
$ {n\choose k}{k\choose j}
 ={n\choose j}{n-j\choose k-j}$ を用いた．
最後の和は $(1-1)^{n-j}$ であり，
$j=n$ なら $1$，それ以外なら $0$ となる．
従って結果は $g_n$ である．
逆方向も同じ有限和を逆順に整理すれば
二項係数行列と符号を付けた行列の積が恒等行列になるため成立する．
\end{proof}
\begin{proposition}[多項式の有限差分]
$R$ を有理数体を含む可換環とし，
$\Delta f(x)=f(x+1)-f(x)$ と定める．
任意の多項式 $f(x)\in R[x]$ の次数が $d$ 以下なら
$\Delta^{d+1}f=0$ である．
\end{proposition}
\begin{proof}
$x^k$ に対して二項定理から
$(x+1)^k-x^k$ の次数は高々 $k-1$ である．
$\Delta$ は $R$-線形なので，
一般の次数 $d$ 以下の多項式に作用させると
次数が少なくとも一つ下がる．
よって $d+1$ 回繰り返すと零多項式になる．
\end{proof}

\section{有限半順序集合の接続代数}
$P$ を有限半順序集合，$R$ を単位元を持つ可換環とする．
\begin{definition}
$\mathcal I(P)=\{(x,y)\in P^2:x\le y\}$ と置く．
$\mathcal I(P)$ 上の $R$ 値関数全体を $I(P;R)$ と書く．
その二元 $f,g$ について畳み込み積を
\[
 (f*g)(x,y)=\sum_{x\le z\le y}f(x,z)g(z,y)
\]
によって定める．
\end{definition}
\begin{lemma}[結合律と単位元]\label{lem:incidence}
畳み込み積は結合的である．さらに
\[
 \delta(x,y)=
 \begin{cases}1&(x=y),\\0&(x<y)\end{cases}
\]
は両側単位元となる．
\end{lemma}
\begin{proof}
$((f*g)*h)(x,y)$ を展開すると
\[
 \sum_{x\le u\le z\le y}f(x,u)g(u,z)h(z,y)
\]
となる．$f*(g*h)$ の展開も同じ添字集合上の和なので，
有限和の順序変更により結合律が成り立つ．
$\delta*g$ の和で非零なのは $z=x$ の項のみであり，
$g*\delta$ では $z=y$ の項のみである．
よって両方とも $g$ に等しい．
\end{proof}
\begin{definition}
全ての $x\le y$ で $\zeta(x,y)=1$ と置き，
$\eta=\zeta-\delta$ とする．
\end{definition}
\begin{lemma}[狭義上三角関数の冪零性]\label{lem:nil}
$m=|P|$ とすると $\eta^{*m}=0$ である．
\end{lemma}
\begin{proof}
$\eta^{*k}(x,y)$ は
$x=x_0<x_1<\cdots<x_k=y$ を満たす
狭義増加列の個数である．
これは畳み込みの定義から $k$ に関する帰納法で従う．
$k\ge m$ なら異なる $k+1>m$ 個の元を必要とし，
そのような列は存在しない．従って零である．
\end{proof}
\begin{theorem}[M\"obius関数の存在と一意性]\label{thm:mobius}
\[
 \mu=\sum_{k=0}^{m-1}(-1)^k\eta^{*k}
 \tag{1}
\]
と定めると
$\mu*\zeta=\zeta*\mu=\delta$ となる．
この性質を持つ $\mu$ は一意であり，
$\mu(x,x)=1$，$x<y$ のとき
\[
 \mu(x,y)=-\sum_{x\le z<y}\mu(x,z)
 \tag{2}
\]
を満たす．
\end{theorem}
\begin{proof}
$\zeta=\delta+\eta$ と結合律によって
\[
 (\delta+\eta)*
 \sum_{k=0}^{m-1}(-1)^k\eta^{*k}
 =\delta+(-1)^{m-1}\eta^{*m}=\delta.
\]
右側から掛ける場合も同じ計算である．
両側逆元は一意である．
$\mu*\zeta=\delta$ の $(x,y)$ 成分は
$\sum_{x\le z\le y}\mu(x,z)=\delta(x,y)$ なので，
$x=y$ で値が $1$，$x<y$ で式 (2) を得る．
\end{proof}
\begin{remark}
式 (1) と補題\ref{lem:nil}から
\[
 \mu(x,y)=\sum_{k\ge0}(-1)^k
 \#\{x=x_0<x_1<\cdots<x_k=y\}
\]
が従う．各区間 $[x,y]$ が有限な
局所有限半順序集合でも，区間ごとに同じ計算ができる．
\end{remark}

\section{M\"obius反転}
\begin{theorem}[半順序集合上の反転定理]\label{thm:posetinv}
有限半順序集合 $P$ 上の関数 $f,g:P\to R$ について
\[
 f(y)=\sum_{x\le y}g(x)\qquad(y\in P)
 \tag{3}
\]
であることと
\[
 g(y)=\sum_{x\le y}\mu(x,y)f(x)\qquad(y\in P)
 \tag{4}
\]
であることは同値である．
\end{theorem}
\begin{proof}
式 (3) を式 (4) に代入すると
\[
 \begin{aligned}
 \sum_{x\le y}\mu(x,y)f(x)
 &=\sum_{x\le y}\mu(x,y)\sum_{z\le x}g(z)\\
 &=\sum_{z\le y}g(z)\sum_{z\le x\le y}\mu(x,y)\\
 &=\sum_{z\le y}g(z)(\zeta*\mu)(z,y)=g(y).
 \end{aligned}
\]
最後は定理\ref{thm:mobius}による．
逆向きには式 (4) を (3) に代入し，
有限和を交換して $\mu*\zeta=\delta$ を用いる．
\end{proof}
\begin{example}[Boolean束]
有限集合 $E$ の部分集合全体を包含関係で順序付ける．
$A\subset B$ に対し
\[
 \mu(A,B)=(-1)^{|B\setminus A|}
\]
となる．
\end{example}
\begin{proof}
$x=y$ の場合には1である．
$A\subsetneq B$ なら
\[
 \sum_{A\subset C\subset B}(-1)^{|C\setminus A|}
 =\sum_{j=0}^{|B\setminus A|}
 {\,|B\setminus A|\choose j}(-1)^j
 =0.
\]
従ってこの関数と $\zeta$ との畳み込みは $\delta$ であり，
逆元の一意性から結論を得る．
\end{proof}
\begin{corollary}[部分集合の包含に関する反転]
有限集合 $E$ の部分集合上の関数で
$f(B)=\sum_{A\subset B}g(A)$ なら
\[
 g(B)=\sum_{A\subset B}
 (-1)^{|B\setminus A|}f(A).
\]
\end{corollary}
\begin{proof}
Boolean束のM\"obius関数を
定理\ref{thm:posetinv}に代入する．
\end{proof}
\section{整除関係と算術関数}
\begin{definition}
正整数 $n$ に対して
\[
 \mu_{\mathrm{ar}}(n)=
 \begin{cases}
 1&(n=1),\\
 (-1)^r&(n\text{ が相異なる }r\text{ 個の素数の積}),\\
 0&(n\text{ がある素数の平方で割り切れる})
 \end{cases}
\]
と定める．
\end{definition}
\begin{lemma}\label{lem:arithmu}
\[
 \sum_{d\mid n}\mu_{\mathrm{ar}}(d)
 =\begin{cases}1&(n=1),\\0&(n>1).\end{cases}
\]
\end{lemma}
\begin{proof}
$n=1$ の場合は定義による．
$n>1$ の相異なる素因数を $p_1,\dots,p_r$ とする．
和に寄与する約数は素因数の重複のない積に限られるので
\[
 \sum_{d\mid n}\mu_{\mathrm{ar}}(d)
 =\sum_{J\subset\{1,\dots,r\}}(-1)^{|J|}
 =(1-1)^r=0.
\]
\end{proof}
\begin{theorem}[整除順序のM\"obius反転]\label{thm:arithinv}
$F,G:\mathbb N_{\ge1}\to R$ に対して
\[
 F(n)=\sum_{d\mid n}G(d)
\quad\Longleftrightarrow\quad
 G(n)=\sum_{d\mid n}\mu_{\mathrm{ar}}(n/d)F(d).
\]
\end{theorem}
\begin{proof}
左辺を右辺に代入して有限和を交換すると
\[
 \begin{aligned}
 \sum_{d\mid n}\mu_{\mathrm{ar}}(n/d)F(d)
 &=\sum_{e\mid n}G(e)
   \sum_{e\mid d\mid n}\mu_{\mathrm{ar}}(n/d)\\
 &=\sum_{e\mid n}G(e)\sum_{a\mid n/e}\mu_{\mathrm{ar}}(a)
 =G(n)
 \end{aligned}
\]
である．最後は補題\ref{lem:arithmu}による．
逆方向も同じ計算を逆向きに行えばよい．
\end{proof}
\begin{proposition}[Eulerの$\varphi$関数]
$\varphi(n)=|\{1\le k\le n:\gcd(k,n)=1\}|$ と置くと
\[
 \sum_{d\mid n}\varphi(d)=n,\qquad
 \varphi(n)=n\sum_{d\mid n}{\mu_{\mathrm{ar}}(d)\over d}
 =n\prod_{p\mid n}\left(1-{1\over p}\right).
\]
\end{proposition}
\begin{proof}
$1,\dots,n$ の各 $k$ を
$d=n/\gcd(k,n)$ によって分類する．
そのような $k$ は $\varphi(d)$ 個であるため第一式を得る．
反転定理に $F(n)=n,G(n)=\varphi(n)$ を適用すると
$\varphi(n)=\sum_{d\mid n}\mu_{\mathrm{ar}}(n/d)d$．
$d$ を $n/d$ へ置換して第二式の和を得る．
最後の有限積は平方因子を含まない約数だけが寄与することから従う．
\end{proof}
\begin{remark}
包除原理と二項反転はBoolean束上の反転の現れである．
反転公式の適用では，上方集合と下方集合のどちらを
足し合わせているかを明示することが重要である．
\end{remark}

\section{到達課題の完全解答}
\noindent\textbf{全射の個数を包除原理から求める．}
\begin{proposition}
5元集合から3元集合への全射は150個ある．
\end{proposition}
\begin{proof}
全ての写像は $3^5$ 個である．
終域の指定された一元を像から除く写像は
$2^5$ 個であり，これを三元について引く．
二つの指定された元を像から除く写像は1個で，
三つの二元部分集合に対応する．
三元すべてを像から除く写像は存在しない．
包除原理により
\[
3^5-\binom31\,2^5+\binom32\,1^5
=243-96+3=150
\]
となる．
\end{proof}
\begin{proposition}[一般の全射計数]
$n$ 元集合から $m$ 元集合への全射数は
\[
\sum_{j=0}^m(-1)^j\binom mj(m-j)^n
\]
である．ただし $n\ge1$ とする．
\end{proposition}
\begin{proof}
終域の各点 $a$ を像から除く写像の集合を
$E_a$ と置く．$j$ 個の点を除くと残りの
$m-j$ 点に各入力を送るので写像数は $(m-j)^n$．
全射とはどの $E_a$ にも属さない写像であるから
包除原理から表示式を得る．
\end{proof}
\noindent\textbf{到達確認．}
3元集合から2元集合への全射数を求めよ．
\par\noindent\textbf{解答．}
$2^3-2\cdot1^3=6$ である．
\end{document}
